% ============================================================
% 第 11 章 电机的发热与温升限值（教材第三篇，印刷页 116--121）
% 节名：第一节 电机的温升和温升限值；第二节 电机的发热与冷却过程；第三节 电机的工作制；第四节 电机的冷却方式
% ============================================================
\chapter{电机的发热与温升限值}

前面各章只关心电磁关系，本章转向\textbf{热}的问题。电机的容量并不完全由电磁参数决定，往往受限于绝缘材料能承受的温度：损耗产生热量，热量使绕组温度升高，温度超过绝缘的耐热极限，绝缘就会加速老化甚至击穿。因此"温升限值"是电机设计和使用中一条不可逾越的红线。

\section{电机的温升和温升限值}

\subsection{温升的定义}
电机运行中产生的各种损耗（铜耗、铁耗、机械损耗）最终都变成热能，一部分使电机温度升高，一部分散入周围介质。电机某部分的温度 $\theta$ 与周围冷却介质温度 $\theta_0$ 之差，称为该部分的\textbf{温升}：
\begin{equation}
\tau=\theta-\theta_0 .
\end{equation}
温升的单位是 K（开尔文）或 ℃（作为温差时两者数值相同）。工程上关心温升而不是绝对温度，因为散热能力取决于温差，而绝缘老化速度也主要取决于温度高出环境多少。

\subsection{绝缘材料的耐热等级}
电机绕组绝缘按耐热能力分为若干等级，各级有其"最高允许温度"。常用等级如下（按 GB 规定的耐热分级）。

\begin{center}
\begin{tabular}{@{}llll@{}}
\toprule
耐热等级 & 最高允许温度 & 主要材料 & 典型应用 \\
\midrule
A & 105\,℃ & 浸渍过的棉纱、纸、丝绸等 & 老式小电机 \\
E & 120\,℃ & 环氧树脂、聚酯薄膜等 & 中小型电机 \\
B & 130\,℃ & 云母、玻璃丝、石棉与漆结合 & 常用等级 \\
F & 155\,℃ & 云母、玻璃丝与耐热漆 & 现代电机主流 \\
H & 180\,℃ & 硅有机树脂、云母、玻璃丝 & 高温、起重、冶金 \\
C & $>180$\,℃ & 云母、玻璃、陶瓷等无机材料 & 特殊场合 \\
\bottomrule
\end{tabular}
\end{center}

\subsection{温升限值}
\textbf{温升限值}是标准规定的、电机各部分温升不得超过的最大值。它由绝缘的极限温度扣除规定的冷却介质最高温度（我国规定环境空气温度最高为 40\,℃）以及绕组内部的\textbf{热点温差}确定。

例如 B 级绝缘、最高允许温度 130\,℃，若环境温度按 40\,℃ 计，则平均温升限值约为 80\,K，再考虑绕组最热点比平均温度高若干度（热点温差），标准给出的温升限值略低于这个差。这就是为什么铭牌上写"温升 80\,K"而不是"温度 130\,℃"。

\begin{center}
\begin{tabular}{@{}lll@{}}
\toprule
项目 & 含义 & 数值举例（B 级） \\
\midrule
最高允许温度 & 绝缘长期工作不加速老化的上限 & 130\,℃ \\
环境温度上限 & 标准规定的冷却介质最高温度 & 40\,℃ \\
温升限值 & 标准规定的允许温升 & 约 80\,K（扣除热点温差后） \\
\bottomrule
\end{tabular}
\end{center}

\corepoint{三条要分清：\textbf{温度}是绝对量，\textbf{温升}是温差，\textbf{温升限值}是标准规定值。考试常问"环境温度 40\,℃ 时 B 级绝缘电机绕组允许的最高温度是多少"，答案不是简单的 130\,℃ 也不是 120\,℃，而要按"温升限值 + 环境温度"算；反过来问"温升限值"，要用"极限温度 $-$ 环境温度 $-$ 热点温差"。}

\section{电机的发热与冷却过程}

\subsection{发热过程}
设电机某部分在单位时间内产生的热量为 $Q$，其热容量为 $C$（使该部分温度升高 1\,K 所需热量），散热系数为 $A$（单位温升、单位时间散出的热量）。由能量守恒，发热与散热之差等于热容量的增温：
\begin{equation}
Q\,\mathrm{d}t=C\,\mathrm{d}\tau+A\tau\,\mathrm{d}t ,
\end{equation}
解得温升随时间的变化规律为
\begin{equation}
\tau=\tau_\infty\left(1-\mathrm{e}^{-t/T}\right),\qquad
\tau_\infty=\frac{Q}{A},\qquad
T=\frac{C}{A}.
\end{equation}

\keyformula{\tau=\tau_\infty\left(1-\mathrm{e}^{-t/T}\right),\qquad T=\frac{C}{A}}

其中 $\tau_\infty$ 为\textbf{稳态温升}（发热与散热达到平衡时的温升），$T$ 为\textbf{发热时间常数}。温升按指数曲线上升：起始上升快，逐渐变慢，理论上 $t\to\infty$ 才达到 $\tau_\infty$，工程上认为 $t=(3\sim4)T$ 时已基本稳定（达 95\%～98\%）。

\subsection{冷却过程}
电机从运行状态切除电源（或负载减小）后，$Q=0$，温升按指数曲线下降：
\begin{equation}
\tau=\tau_0\mathrm{e}^{-t/T},
\end{equation}
其中 $\tau_0$ 为冷却开始时的温升。冷却时散热条件往往变差（如风扇停转），故冷却时间常数一般大于发热时间常数。

\subsection{时间常数的工程意义}
\begin{itemize}
  \item 小型电机发热时间常数约为十几分钟到一小时；
  \item 中大型电机可达数小时。
\end{itemize}
时间常数越长，说明电机"热惯性"越大：短时过载不会立即造成危险，但过载持续时间一长就危险。这正是"短时工作制"与"连续工作制"容量可以不同的原因。

\section{电机的工作制}

电机的工作制（工作方式）指电机承受负载的情况，它决定了电机的发热与冷却循环，也决定了额定功率的标定条件。按 GB 规定，主要有三类。

\subsection{连续工作制（S1）}
电机在恒定负载下连续运行，温升达到稳态值。绝大多数工业用电动机（水泵、风机、压缩机）属此类。铭牌上的额定功率就是在连续工作制、规定的环境温度和冷却条件下，能长期输出而不超过温升限值的功率。

\subsection{短时工作制（S2）}
电机在恒定负载下运行一段有限时间（尚未达到稳态温升），随后断电停转一段时间，温升降到接近环境温度。如闸门启闭机、某些机床辅助传动。

\subsection{断续周期工作制（S3）}
电机按"运行—停歇—运行"周期性循环，每次运行时间不足以达到稳态温升，停歇时间也不足以冷却到环境温度。用\textbf{负载持续率}（暂载率）$FC$ 表征：
\begin{equation}
FC=\frac{t_r}{t_r+t_0}\times100\%,
\end{equation}
其中 $t_r$ 为运行时间、$t_0$ 为停歇时间，一个周期一般不超过 10\,min。标准负载持续率有 15\%、25\%、40\%、60\% 等。起重机、电梯、轧机等属此类。

\corepoint{同一台电机，工作制不同，允许输出的功率不同。短时或断续工作的电机，因温升来不及达到稳态，可在短时间内输出比连续额定功率更大的功率；反之，把连续工作制的电机用于短时工作制是"浪费"，把短时工作制的电机用于连续运行则会过热烧毁。看到铭牌功率，先问"是哪种工作制下的额定值"。}

\section{电机的冷却方式}

\subsection{按冷却介质分}
\begin{itemize}
  \item \textbf{空气冷却}：最常见，分自扇冷（转子自带风扇）与他扇冷（独立风机）。
  \item \textbf{氢气冷却}：用于大型汽轮发电机。氢气的导热系数与比热容比空气大，且密度小，可显著提高散热能力并降低通风损耗；但必须严格密封、防止爆炸。
  \item \textbf{水冷却}：水的比热容和导热系数远大于空气，用于大型发电机定子绕组（水内冷）及大功率变频电机。
\end{itemize}

\subsection{按防护与通风结构分}
\begin{itemize}
  \item \textbf{开启式}：与外界空气直接相通，散热好但易进尘、受潮。
  \item \textbf{防护式}：有网罩或挡板，防止固体异物进入。
  \item \textbf{封闭式}：机壳封闭，靠机壳表面散热或内部循环冷却，防尘防潮好，但散热能力较差，同容量体积较大。
\end{itemize}

\subsection{通风路径}
按冷却气体在电机内部的流动路径，可分为\textbf{径向通风}（气体沿径向穿过铁芯与绕组）、\textbf{轴向通风}（沿轴向流过）和\textbf{混合通风}。大型电机常在铁芯中留出径向通风沟，把铁芯沿轴向分成若干段，以增大散热面积。

\begin{center}
\begin{tabular}{@{}p{2.4cm}p{4.2cm}p{5.2cm}@{}}
\toprule
项目 & 内容 & 关键结论 \\
\midrule
温升 & $\tau=\theta-\theta_0$ & 绝缘老化取决于温度，故标准规定温升限值 \\
绝缘等级 & A/E/B/F/H/C & 最高允许温度 105/120/130/155/180/$>180$\,℃ \\
发热过程 & $\tau=\tau_\infty(1-\mathrm{e}^{-t/T})$ & $3\sim4$ 个时间常数后基本稳定 \\
冷却过程 & $\tau=\tau_0\mathrm{e}^{-t/T}$ & 冷却时间常数通常大于发热时间常数 \\
工作制 & S1 连续 / S2 短时 / S3 断续 & 断续用负载持续率 $FC$ 表征 \\
冷却方式 & 空气、氢、水；开启、防护、封闭 & 容量越大越依赖强化冷却 \\
\bottomrule
\end{tabular}
\end{center}
